
\documentclass[12pt]{article}
\usepackage{amsmath, amssymb}
\usepackage{geometry}
\geometry{margin=1in}
\title{SAT.O4 – Topological Falsifiability of Particle Classes}
\author{SATO.4D Coherence Project}
\date{June 2025}

\begin{document}
\maketitle

\section*{1. Objective}

To establish falsifiable predictions from the SAT framework based on the allowed topological configurations of filament bundles in 4D.

\section*{2. Core Claim}

Only the following topological filament configurations yield physically stable bound states:
\begin{itemize}
    \item \( n = 1 \): Unbound filament (e.g. neutrino analog)
    \item \( n = 2 \): Linked pair (Hopf link, meson class)
    \item \( n = 3 \): Triple-linked structure (Borromean ring, baryon class)
\end{itemize}

Any observed stable particle with \( n \geq 4 \) would falsify the SAT framework.

\section*{3. Binding Stability}

Stability is defined via topological invariance under smooth 4D deformations:
\[
\frac{\delta Q}{\delta \Sigma_t} = 0
\]
where \( Q \) is a linking/winding/writhing invariant defined for the bundle, and \( \Sigma_t \) is the resolving surface.

\section*{4. Phase Locking and Integer Classes}

Filaments bind only under:
\[
\phi_i - \phi_j = \frac{2\pi k}{n}, \quad k \in \mathbb{Z}
\]
Stable phase-locked configurations exist only for \( n \leq 3 \). Larger \( n \) leads to internal instability or kinematic decay.

\section*{5. Predictive Bound}

SAT predicts no stable filament composites beyond:
\[
n_{\text{max}} = 3
\]

This rule holds across all energy scales and provides a strong falsifiability constraint.

\section*{6. Implications for New Physics}

\begin{itemize}
    \item Any observation of bound states with \( n > 3 \) (e.g., tetraquarks, pentaquarks, exotic hadrons) must be interpreted as unstable resonances or projections of lower-n bundles.
    \item Confirmed existence of topologically stable tetra-bundles would refute SAT.O4.
\end{itemize}

\section*{7. Cross-Module Consistency}

\begin{itemize}
    \item \textbf{O3}: Gauge group classes terminate at SU(3) due to this topological constraint.
    \item \textbf{O5}: Couplings emerge from linking densities capped at triplet configurations.
    \item \textbf{O8}: Mass suppression Q scaling terminates at \( Q = 3 \).
\end{itemize}

\section*{8. Frame Declaration}

This falsifiability rule is established in \textbf{Interpretive Mode 1}. All limits derive from bundle topology; no symmetry or mass condition is imposed externally.

\end{document}
